\chapter{Conclusion}
\label{chap:conclusion}

This chapter should synthesize the thesis as a whole rather than repeat the final section of each paper.

\section{Summary}

\TODO{Summarize the thesis-level contributions across all included chapters.}

\section{Limitations}

\TODO{Discuss important limitations, open questions, or boundary conditions.}

The PrivPGD results in Chapter~\ref{chap:privpgd} depend on the regime in which the method is used. Its main empirical advantage is clearest at the default privacy level and for large marginal workloads where geometry-aware optimization is useful. At smaller privacy budgets, the privatized marginals become noisier, and PGM-based methods or Private GSD can be more competitive, especially on query- and covariance-based metrics and on smaller or higher-dimensional datasets. Thus, PrivPGD should be read as a flexible generation strategy rather than a uniformly dominant synthesizer across all privacy regimes, datasets, and evaluation metrics.

PrivPGD also relies on a chosen embedding of the tabular domain. The default ordered embedding is natural for continuous or ordered features, but categorical variables without a meaningful order may require different embeddings. For example, race categories might be better represented by equidistant centers, while geographic variables may call for embeddings that preserve spatial distances. These choices affect which discrepancies the Wasserstein objective treats as large or small, and therefore form part of the modeling assumptions behind the method.

CP-Fuse, introduced in Chapter~\ref{chap:cpfuse}, relies on a different kind of structural assumption: the copyrighted material must be sufficiently separable across the base models. In practice, this may be difficult to guarantee. Copyrighted passages can appear in duplicated, quoted, reformatted, or near-verbatim forms across a training corpus, and identifying all such replicas before partitioning the data is itself a hard curation problem. The overlap experiments in Appendix~\ref{app:cpfuse-separability} suggest that moderate overlap can still leave useful protection, but the protection deteriorates sharply when the two models memorize substantially overlapping protected material.

The method should therefore be understood as a technical mitigation for regurgitation, not as a complete legal solution to copyright risk. The experiments measure exact and approximate copying in controlled text and code settings, including deliberately overfitted models and prefix-prompt extraction attacks. They do not by themselves determine when a generation is legally infringing, nor do they fully cover deployment settings involving larger models, broader corpora, licensing restrictions, or adversaries with stronger extraction strategies.

The framework for quantifying hidden confounding in Chapter~\ref{chap:confounding-test} should be interpreted in light of its assumptions. It relies on transportability of the conditional average treatment effect between the randomized trial and the observational study. If this assumption is violated, then a discrepancy between randomized and observational evidence may be attributed to hidden confounding even when it instead reflects genuine effect heterogeneity across study populations.

The resulting lower bound is also optimistic outside the common support of the two studies. The procedure can only use randomized evidence for target populations that are sufficiently supported by the trial. Hidden confounding may therefore be larger in regions of the observational covariate space not covered by the randomized study. Finally, the tests are designed to detect hidden confounding that biases the target average treatment effect. Confounding effects that cancel out on average may remain invisible, even if they matter for more local or subgroup-level causal quantities. Relatedly, the marginal sensitivity model can be conservative when the hidden confounder is only weakly related to the potential outcomes, leading to lower bounds that are valid but loose.

\section{Future Work}

\TODO{Outline the most promising next steps that follow from the thesis.}

For PrivPGD, future work could develop regularization penalties that provide stronger inductive bias, analogous to the maximum-entropy bias of PGM methods, while preserving the flexibility of particle optimization. A second direction is to design embeddings and differentiable constraints tailored to specific application domains, such as categorical features without natural order or geographic variables with meaningful metric structure. Finally, practical deployments would benefit from guidance on when to prefer PrivPGD over PGM-based or query-based alternatives as privacy budgets, dataset sizes, dimensionality, and target metrics vary.

For CP-Fuse, the most direct next step is to understand the method under partial separability. This includes theoretical bounds on the longest copied segment as overlap between protected partitions increases, as well as practical data-curation procedures for finding replicas and assigning protected material to partitions. Another direction is to evaluate CP-Fuse in larger-scale deployment settings with genuine copyright-sensitive corpora, such as books, songs, or licensed codebases. Finally, the model-fusion perspective suggests extensions to more than two base models and to hybrid systems that combine CP-Fuse with training-time safeguards, retrieval controls, or legal-review workflows.

For the confounding-strength lower bound, these limitations point to several extensions. One direction is to develop richer sensitivity models that account not only for treatment-assignment distortion, but also for how informative hidden confounders are about potential outcomes. Another is to extend the approach to settings with multiple observational datasets but no randomized trial. A third direction is to move from diagnosis to intervention: once a study is flagged as substantially confounded, it would be valuable to identify which covariates, design changes, or target-population restrictions would most effectively mitigate the problem.
